\documentclass[11pt]{article}
\usepackage[margin=1in]{geometry}
\usepackage{amsmath,amssymb,bm,booktabs,array}
\usepackage[T1]{fontenc}
\usepackage{lmodern}
\setlength{\parindent}{0pt}
\setlength{\parskip}{0.65em}
\newcommand{\R}{\mathbb{R}}
\begin{document}

\begin{center}
{\LARGE\bfseries MIT VNAV 2023 -- Lab 3 Individual Solutions}\\[0.6em]
{\large Wenqian Zhang \quad Student ID: 24020036087}
\end{center}
\hrule

\section*{Deliverable 1: Quaternion API}
The concise code answers are supplied in \texttt{deliverable\_1.txt}. The key
conversions are \texttt{tf2::fromMsg}, \texttt{tf2::toMsg},
\texttt{tf2::getYaw}, and \texttt{tf::quaternionMsgToEigen}. The scalar
component of a \texttt{tf2::Quaternion} is returned by \texttt{getW()}.

\section*{Deliverable 2A: Rank of the Quadrotor Allocation Matrix}
Let each rotor generate thrust \(c_f\omega_i^2\bm e_3\) and reaction torque
\(s_i c_d\omega_i^2\bm e_3\), where \(s_i\in\{+1,-1\}\) records the spin
direction. For rotor position \(\bm\rho_i=(x_i,y_i,0)^T\), its contribution
to the reduced wrench \((f_z,\tau_x,\tau_y,\tau_z)^T\) is
\[
\begin{bmatrix}c_f&c_fy_i&-c_fx_i&s_ic_d\end{bmatrix}^T\omega_i^2.
\]

Use the X configuration in the figure, with arm length \(l\), rotor order
\(1,2,3,4\), and positions
\[
\bm\rho_1=\frac l{\sqrt2}(1,1,0)^T,\quad
\bm\rho_2=\frac l{\sqrt2}(-1,1,0)^T,
\]
\[
\bm\rho_3=\frac l{\sqrt2}(-1,-1,0)^T,\quad
\bm\rho_4=\frac l{\sqrt2}(1,-1,0)^T.
\]
Set \(a=c_fl/\sqrt2\) and \(d=c_d\). With the normal alternating spin
directions \((+,-,+,-)\), the nonzero part of the allocation matrix is
\[
\bar F_a=
\begin{bmatrix}
c_f&c_f&c_f&c_f\\
a&a&-a&-a\\
-a&a&a&-a\\
d&-d&d&-d
\end{bmatrix}.
\]
The full six-dimensional wrench matrix simply adds two zero force rows, so it
has the same column rank as \(\bar F_a\). For \(c_f,a,d\ne0\), the four rows
above are independent; hence
\[
\operatorname{rank}(F_a)=\operatorname{rank}(\bar F_a)=4.
\]
Thus the four squared rotor speeds independently control total thrust and the
three body torques.

In configuration (b), rotors 3 and 4 reverse direction, so the yaw-sign pattern
becomes \((+,-,-,+)\):
\[
\bar F_b=
\begin{bmatrix}
c_f&c_f&c_f&c_f\\
a&a&-a&-a\\
-a&a&a&-a\\
d&-d&-d&d
\end{bmatrix}.
\]
Its last row is \(-d/a\) times its third row. The remaining three rows are
independent, so
\[
\operatorname{rank}(F_b)=\operatorname{rank}(\bar F_b)=3.
\]
The vehicle therefore loses one independent control direction: the yaw moment
is coupled to the pitch moment. This is why the modified vehicle cannot command
an arbitrary four-component thrust/torque input.

\section*{Deliverable 2B: Aerodynamic-Drag Compensation}
Let \({}^{b}\bm v=({v_x},{v_y},{v_z})^T\) be body-frame velocity and let
\(D=\operatorname{diag}(0.1,0.1,0.2)\) contain the given positive drag
coefficients.
The componentwise quadratic model can be written
\[
{}^{b}\bm F_{\mathrm{drag}}
=-D\begin{bmatrix}v_x|v_x|\\v_y|v_y|\\v_z|v_z|\end{bmatrix}.
\]
If \({}^{w}R_b\) maps body coordinates to world coordinates, then
\[
{}^{w}\bm F_{\mathrm{drag}}
={}^{w}R_b\,{}^{b}\bm F_{\mathrm{drag}}.
\]

Let the nominal geometric-controller force be, for example,
\[
\bm F_{\mathrm{nom}}
=-K_x\bm e_x-K_v\bm e_v+mg\bm e_3+m\bm a_d.
\]
Since the physical dynamics contain the additional drag force, use the
feed-forward-compensated command
\[
\boxed{\bm F_{\mathrm{cmd}}
=\bm F_{\mathrm{nom}}-{}^{w}\bm F_{\mathrm{drag}}.}
\]
The desired thrust direction is then
\(\bm b_{3d}=\bm F_{\mathrm{cmd}}/\|\bm F_{\mathrm{cmd}}\|\), and the scalar
collective thrust is obtained by projecting the command onto the current body
thrust axis (with the sign adjusted if the course uses thrust along
\(-\bm e_3\)). This feed-forward term cancels the modeled aerodynamic force,
reduces the steady tracking error at speed, and leaves the feedback terms to
correct modeling errors and disturbances.

\section*{Submission Scope}
Only the individual deliverables are included here. The Lab 3 controller is a
team deliverable and belongs in the team repository, not this personal one.

\end{document}

